\section{Introduction}
\label{sec:intro}

How much of a solid dissolves in a given liquid decides whether a drug is formulated as a tablet or
an injection, which solvent a crystallisation runs in, and how a product is purified. Each value is
measured one solute, one solvent and one temperature at a time, so most combinations of industrial
interest have never been measured, and a laboratory that needs one is choosing between running an
experiment and trusting a prediction.

Predictors divide by how much physical structure they impose. A direct model regresses solubility
from descriptors or from a representation learned off the molecular graph, and on extrapolation to
unseen solutes the direct models lead; FastSolv \cite{osanloo2024} reports the lowest errors
on that task. What such a model returns is a single number, silent about which part of the
chemistry produced it. A physics-informed model instead predicts the quantities the physical
picture names: an activity coefficient for the solute--solvent interaction, a melting temperature
and an enthalpy of fusion for the crystal. A fixed thermodynamic relation then combines them.
SolProp \cite{vermeire2022solprop} assembles such a cycle from separately learned components. The fixed relation costs accuracy and buys three things. A prediction made
through it can be taken apart, because each intermediate is a quantity chemists tabulate on its
own. Each intermediate can be supervised against reference data that a direct regressor has no way
to use. And the relation is itself a prior, asserting for free what a purely statistical model
would have to learn from data it does not have, which is the situation whenever the molecule in
question is unlike anything in the training set.

Among those intermediates the $\sigma$-profile is the one most often learned. In the
conductor-like screening model for real solvents and in its segment-activity reformulation
COSMO-SAC \cite{lin2002}, a molecule enters the thermodynamics only through the distribution of
screening-charge density over its cavity surface: how much of that surface carries charge density
of each value. The distribution comes from a quantum-chemical calculation with a continuum
solvation treatment and is tabulated once per molecule \cite{mullins2006}, after which a
closed-form segment-activity expression consumes it and knows nothing else about the compound.
That makes it an unusually clean target for learning. It is a fixed-length vector, it is tabulated
for thousands of molecules, and supplying it from structure removes the only expensive step in the
pipeline. Group-contribution schemes supplied it first \cite{mu2007gc,mu2009gc}, and learned
predictors followed, taking the profile from a line notation or a molecular graph
\cite{kang2022surrogate,chen2021transformer,zhang2022mpnn}. How far the tabulated reference itself
moves with the quantum chemistry and solvation treatment behind it is a separate and live question
\cite{salih2025sigma}. An older line regresses an \emph{apparent} profile from phase-equilibrium
data instead \cite{islam2015apparent}, which is the pre-neural form of the effective profile
measured here.

All of this rests on a premise rarely written down, because it reads as though it needs no writing.
If the quantity in the slot is a $\sigma$-profile, then a network that fills the slot has learned a
$\sigma$-profile. The usual consequences follow: it can be read as chemistry, and it can be
compared against tabulation. The strongest consequence is that a better value can be looked up and
handed back to the model without retraining, and the prediction should improve. The premise is what makes the
intermediate worth having, since it is the difference between a physical model and an elaborate
parameterisation. The recent literature states it in close to those words. The physics-embedded
phase-equilibrium model closest to the work reported here argues that its intermediate carries
clear physical significance because it appears directly in the thermodynamic formulation
\cite{yang2025tennet}. That argument is about the equations. Whether the values a network places in
the slot are the physical quantity is a different question, and it takes a different kind of
answer.

There is reason to expect that they are not. In machine learning the same construction is called a
concept bottleneck \cite{koh2020cbm}. A network is required to predict named, human-meaningful
concepts, and the target is predicted from those concepts alone, in the expectation that the
prediction becomes interpretable and that intervening on a concept behaves as its name implies. The
concepts leak. A bottleneck trained end to end carries information its labels do not name, the
downstream model reads it for more than its nominal content, and both the interpretation and the
interventions then fail in ways the names do not predict
\cite{margeloiu2021,mahinpei2021,havasi2022}. A $\sigma$-profile head inside a fixed closure is
that construction with a thermodynamic decoder in place of a learned one, and it differs in one
respect that matters: the concept has a ground truth. A $\sigma$-profile is not a human label but a
computable quantity, so leakage can be measured against the quantity itself rather than inferred
from downstream behaviour. The structural reason for leakage is also sharper here. A solubility
measurement constrains the \emph{composition} of the intermediate with the fixed relation, not the
intermediate itself. Where different intermediates compose to the same target, the data cannot
choose between them, and the combinations it leaves undetermined are where a flexible network is
free to move. This is the familiar predicament of a model whose parameters are weakly constrained
by its observables \cite{gutenkunst2007,transtrum2015,raue2009}, compounded by calibration against
a relation that is itself only approximate \cite{kennedy2001,brynjarsdottir2014}.

The literature has, without saying so, already taken a position on this. In the physics-embedded
model nearest to the present work, the $\sigma$-profile predictor is trained against
quantum-chemical profiles and then held fixed. At the fine-tuning stage every component and layer
is frozen except the output head of the segment-activity network, so no gradient from the
experimental phase-equilibrium data reaches the intermediate \cite{yang2025tennet}. At the
other end of the same axis sit models that learn a thermodynamic potential from equilibrium
compositions with no supervision on the intermediate at all, and models that dispense with the
named intermediate and impose the thermodynamics in the architecture instead
\cite{specht2024hanna}. The axis is a design choice about how far the intermediate is allowed to
move under the target's gradient. It is made differently in nearly every paper and argued for in
none of them. The pipeline studied here sits in the middle by construction: its $\sigma$-profile
head is supervised against the same tabulated reference, and the closure's gradient also reaches
it. That is what makes it a place where the cost of the choice can be read.

This paper asks what a network puts in that slot, on one pipeline and one intermediate, and reports
a battery of measurements rather than a model. The first asks whether solubility determines the
profile at all, and answers through the curvature of the likelihood in the directions the closure
exposes. The second asks whether the head uses the freedom the first one finds, and where on the
charge-density axis it goes. The third asks whether the architecture is responsible, by testing
whether the output layer that cannot represent an exact zero is the reason the learned profile
carries charge density where the reference carries none. The fourth asks whether the resulting
profile still respects the reference state, that is whether a molecule dissolved in itself is
predicted to be ideal. The fifth asks what the misfit costs the closure and where on the grid that
cost is located. The sixth asks whether the intermediate decomposes at all, by replacing parts of
the learned profile with the reference and asking whether the parts add up. The last asks when
supplying a faithful intermediate helps, using a second and better-specified closure for contrast.
Each test is stated with the result that would refute its reading before it is run, because several
of them return negatives, and a negative that is wrong closes a direction instead of correcting
it.

Three things the paper does not claim. It does not report a more accurate solubility model. On the
scaffold-disjoint split used throughout, the physics path trails a matched direct control in the
seed mean, and both sit near a random-forest baseline. Both are also well behind the external state
of the art on extrapolation to new solutes, and the absolute errors are within reach of the
disagreement between independent measurements of the same system. It is not a benchmark against any published
model, because no published model was retrained here; the comparison with prior art is a comparison
of design choices and of what each design leaves unmeasured. And it does not show that learning a
$\sigma$-profile is a poor design. It shows what has to be measured before the learned profile is
read as one.
